\documentclass[10pt,twocolumn]{article}
\usepackage[T1]{fontenc}
\usepackage[utf8]{inputenc}
\usepackage{lmodern}
\usepackage[margin=1.8cm,top=2.4cm,bottom=2cm,columnsep=0.6cm]{geometry}
\usepackage{fancyhdr}
\usepackage{titlesec}
\usepackage{booktabs}
\usepackage{amsmath,amssymb,amsfonts}
\usepackage{microtype}
\usepackage{xcolor}
\usepackage{mdframed}
\usepackage{parskip}
\usepackage{enumitem}
\usepackage{hyperref}

\definecolor{rulered}{RGB}{180,30,30}
\definecolor{boxgray}{RGB}{245,245,245}
\definecolor{keywordgray}{RGB}{100,100,100}

\pagestyle{fancy}
\fancyhf{}
\fancyhead[L]{\textsc{\small Claude Mag}}
\fancyhead[C]{\small\textit{Machine Learning Research Digest}}
\fancyhead[R]{\small 2026-08-17}
\fancyfoot[C]{\small\thepage}
\renewcommand{\headrulewidth}{0.8pt}

\titleformat{\section}{\normalsize\bfseries\color{rulered}}{}{0pt}{}%
  [\vspace{-4pt}\color{rulered}\rule{\columnwidth}{0.4pt}]
\titlespacing{\section}{0pt}{8pt}{4pt}

\hypersetup{colorlinks=true,urlcolor=rulered,linkcolor=black}

\begin{document}
\setlength{\parindent}{0pt}
\setlength{\parskip}{4pt}

{\small\color{keywordgray}\textbf{Keywords:}
  hallucination detection \textbullet{}
  retrieval-augmented generation \textbullet{}
  LLM reliability \textbullet{}
  uncertainty estimation}\par\vspace{4pt}

{\LARGE\bfseries\raggedright
  SIRIN: A Unified Toolkit for Detecting
  Contextual Hallucinations in RAG
  and Memory-Grounded LLM Systems}\par\vspace{4pt}

{\large\itshape\raggedright
  Three Detector Paradigms --- Representation Probing, Uncertainty Estimation,
  and Judge-Style Verification --- Unified Under One Interface and Evaluation Pipeline}\par\vspace{6pt}

{\small\textbf{By} Julia Belikova, Rauf Parchiev, Mikhail Filimonov,
  Konstantin Polev, Andrey Savchenko, Maksim Makarenko}\par\vspace{2pt}

{\small\color{keywordgray}arXiv:2608.00033 \textbullet{} Preprint, August 2026}
\par\vspace{2pt}\noindent{\color{rulered}\rule{\columnwidth}{0.6pt}}\vspace{6pt}

Contextual hallucinations --- fluent, plausible LLM responses that contradict the
provided retrieved evidence --- are the dominant reliability failure in RAG systems.
SIRIN (Semantic Inconsistency Recognition and Inspection Nexus) consolidates three
detector paradigms under one unified toolkit, configuration system, and evaluation
pipeline, supporting both white-box and black-box settings as well as pre-generation
query answerability checks.

\begin{mdframed}[backgroundcolor=boxgray,linecolor=rulered,linewidth=1pt,
    innerleftmargin=6pt,innerrightmargin=6pt,
    innertopmargin=6pt,innerbottommargin=6pt]
\textbf{\small AT A GLANCE}\par\vspace{4pt}
\begin{description}[leftmargin=0pt,labelsep=4pt,itemsep=2pt,parsep=0pt,
    font=\small\bfseries]
  \item[Problem:] Contextual hallucinations in RAG systems are invisible to standard
    quality metrics and require comparing model output to retrieved context.
  \item[Why it matters:] RAG is used in high-stakes applications; hallucination detection
    is safety-critical but existing methods are fragmented across incompatible systems.
  \item[Method:] Unify representation probing, uncertainty estimation, and judge-style
    verification under one interface; add pre-generation query answerability.
  \item[Result:] No single paradigm dominates; combining uncertainty + judge-style
    verification improves precision; answerability filtering reduces hallucination rate.
  \item[Evidence:] Benchmarked across standard RAG hallucination datasets in both
    white-box and black-box settings.
\end{description}
\end{mdframed}

\section{The Big Picture}

Retrieval-augmented generation (RAG) addresses parametric hallucination by grounding
LLM outputs in retrieved documents. But retrieval-grounded generation introduces its
own failure mode: \textbf{contextual hallucination}, where the model generates a response
that is fluent and internally coherent but contradicts or is unsupported by the
specific retrieved evidence. Unlike parametric hallucination (confabulated facts),
contextual hallucination requires detecting inconsistency between model output and
the provided context --- a structured grounding task.

The detection literature is fragmented. Representation probing requires model internals.
Uncertainty estimation can work in grey-box settings. Judge-style verification works
via API but is expensive. No unified system allows fair comparison across paradigms on
the same inputs with the same evaluation protocol --- until SIRIN.

\section{Why Existing Approaches Fall Short}

\textbf{Representation probing} methods train a lightweight classifier on the model's
internal activations to predict hallucination. They are fast and model-specific but
require white-box access and labeled probe training data.

\textbf{Uncertainty estimation} methods measure token-level or span-level prediction
entropy as a proxy for hallucination risk. They are model-agnostic in principle but
require access to logit distributions (grey-box) and can confuse \emph{uncertain}
content with \emph{wrong} content.

\textbf{Judge-style verification} uses a second LLM to verify each claim in the response
against the retrieved context. It is the most model-agnostic and robust but requires
additional LLM calls (cost) and is sensitive to judge model quality.

Each paradigm has a different interface, codebase, and evaluation protocol, making
cross-comparison impossible without SIRIN's unification layer.

\section{Core Mechanism}

SIRIN provides three core components:

\textbf{1. Detector modules} implementing each paradigm:
\begin{itemize}[itemsep=1pt,parsep=0pt]
  \item \textit{Probe detector:} Trains a linear probe on activations from a
    specified layer to classify hallucinated vs.\ grounded spans.
  \item \textit{Uncertainty detector:} Computes semantic entropy, predictive entropy,
    or contrastive uncertainty over the model's token distribution.
  \item \textit{Judge detector:} Prompts a judge LLM to verify each factual claim
    in the response against the retrieved context chunk-by-chunk.
\end{itemize}

\textbf{2. Answerability module:} Before generation, scores whether the retrieved
context contains sufficient information to answer the query. If answerability is below
threshold, the system abstains rather than generates a potentially hallucinated response.

\textbf{3. Unified evaluation pipeline:} Standard datasets (TriviaQA-RAG, NQ-RAG,
HaluEval-RAG), shared metrics (precision, recall, F1 at response and span level),
and a comparison framework that runs multiple detectors on the same inputs.

\section{Technical Formulation}

Let $c$ be the retrieved context, $q$ the query, and $y$ the generated response.
Contextual hallucination detection is a binary classification:
\[
d(y, c) = \begin{cases} 1 & \text{if } y \text{ is inconsistent with } c \\ 0 & \text{otherwise} \end{cases}
\]

\textbf{Probe detector:} Train $f_\phi: \mathbf{h} \to [0,1]$ where
$\mathbf{h} = \mathrm{enc}_\theta(y, c)$ is the model's internal representation:
\[
d_{\mathrm{probe}}(y,c) = \mathbb{1}[f_\phi(\mathbf{h}) > \tau_p]
\]

\textbf{Uncertainty detector:} Estimate semantic entropy as:
\[
H_{\mathrm{sem}} = -\sum_{s \in S} p(s \mid c, q) \log p(s \mid c, q)
\]
where $S$ is the set of semantically distinct responses sampled via beam search or
Monte Carlo sampling. High entropy signals hallucination risk.

\textbf{Judge detector:} For each extracted claim $\hat{y}_i$:
\[
d_{\mathrm{judge}}(\hat{y}_i, c) = \mathrm{LLM}_{\mathrm{judge}}\bigl(``\text{Is }\hat{y}_i\text{ supported by }c\text{?}''\bigr)
\]

\textbf{Answerability:} Pre-generation query answerability score:
\[
a(q, c) = \mathrm{LLM}_{\mathrm{judge}}\bigl(``\text{Does }c\text{ contain enough info to answer }q\text{?}''\bigr)
\]

\section{Learning or Inference Procedure}

\textbf{White-box setting:} Probe detectors require one-time fine-tuning on a labeled
hallucination dataset ($\sim$1k examples). Uncertainty detectors require multiple forward
passes (5--20 samples) per query. No training for the judge detector.

\textbf{Black-box setting:} Only the judge detector is applicable. SIRIN wraps the
judge calls with batching and rate limiting for production deployment.

\section{What the Guarantee Says}

No single detector paradigm dominates across all settings. The key empirical guarantee
is: combining uncertainty estimation + judge-style verification achieves higher
precision than either alone, at moderate additional cost. Answerability filtering
reduces hallucination rate by 15--25\% by preventing generation when context is insufficient.

\section{Experimental Findings}

\begin{itemize}[itemsep=2pt,parsep=0pt]
  \item \textbf{Probe detector:} Best F1 in white-box settings; poor in black-box.
  \item \textbf{Uncertainty detector:} Competitive in grey-box; highest recall but
    lower precision than judge-style.
  \item \textbf{Judge detector:} Most consistent across settings; best in black-box.
  \item \textbf{Combination:} Uncertainty + judge achieves the best precision-recall
    tradeoff across all datasets.
  \item \textbf{Answerability filtering:} Reduces hallucination rate 15--25\%
    by abstaining when context is insufficient.
\end{itemize}

\section{Ablations and Interpretation}

\textbf{Span vs.\ response level:} Span-level detection has higher recall but lower
precision than response-level detection; the interactive UI benefits from span-level
highlighting for analyst workflows.

\textbf{Judge model quality:} Detection performance degrades gracefully as judge
model size decreases; 7B judge models achieve $>$85\% of the performance of
frontier judges at 10$\times$ lower cost.

\textbf{Answerability threshold:} High thresholds ($a > 0.9$) maximize precision
at the cost of high abstention rate; $a > 0.7$ is the practical operating point.

\section*{Reference}

Julia Belikova, Rauf Parchiev, Mikhail Filimonov, Konstantin Polev,
Andrey Savchenko, Maksim Makarenko.
\textbf{SIRIN: A Unified Toolkit for Detecting Contextual Hallucinations
in Retrieval-Augmented and Memory-Grounded LLM Systems}.
arXiv:2608.00033, August 2026.
\url{https://arxiv.org/abs/2608.00033}

\end{document}
